\chapter{Creio na Santa Igreja Católica}
\chaptermark{Creio na Santa Igreja Católica}

Este capítulo aprofunda o terceiro artigo do Símbolo Apostólico, professando a fé na santa Igreja católica. Desdobramos os parágrafos correspondentes do Catecismo da Igreja Católica (CIC 748--945), explorando nomes e origem da Igreja, Povo de Deus, marcas una, santa, católica e apostólica, estrutura hierárquica, leigos e consagrados. Para adultos em processo de iniciação cristã pelo RICA, a Igreja é mãe e mestra, acolhendo no catecumenato rumo aos sacramentos pascais.

\section{Nomes, Imagens e Origem da Igreja no Plano de Deus (\S 748--780)}

A Igreja é ``convocação'' (ekklesia), assembleia litúrgica e universal dos fiéis, herdeira da assembleia do Sinai, chamada por Deus de todos os confins.

Imagens bíblicas revelam seu mistério: Povo de Deus, Corpo de Cristo, Templo do Espírito, Esposa do Cordeiro, Nova Jerusalém, vinhedo, rebanho, casa de Deus.

Plano trinitário: Pai cria para comunhão, Filho convoca em si, Espírito realiza a Igreja como meta da criação, preparada no Antigo Testamento.

Israel elege-se Povo de Deus pela Aliança, figura da nova e eterna em Cristo, reunindo judeus e gentios no Espírito.

A Igreja manifesta-se plenamente em Pentecostes, mas prefigurada desde o início, reação divina ao pecado para restaurar comunhão.

Para adultos a partir dos 20 anos no RICA, essas imagens inspiram o rito de acolhida, unindo catecúmenos à assembleia eucarística como Corpo vivo de Cristo.

\section{O Povo de Deus e suas Características (\S 781--810)}

O Povo de Deus distingue-se: escolhido, não por nascimento carnal mas por Batismo, cabeça em Cristo Rei, lei do amor, missão de sal e luz.

Participa dos ofícios sacerdotal, profético e real de Cristo: santificar, anunciar, reinar com Ele.

Hierarquia (bispos sucessores apostólicos), leigos e consagrados colaboram na missão única da Igreja.

Universal: aberta a todos, peregrina rumo ao Reino, semeando unidade e salvação.

Hierarquia serve unidade: apóstolos e sucessores ensinam, santificam, governam em nome de Cristo.

Na catequese para adultos não iniciados, compreender o Povo de Deus no tempo de catecumenato fortalece identidade batismal, preparando para integração na paróquia como membros ativos.

\section{A Igreja é Una e Santa (\S 811--829)}

A Igreja é una, santa, católica e apostólica, marcas dadas por Cristo via Espírito, professadas no Credo.

Unidade trinitária: fonte no Pai, fundador Cristo, alma o Espírito; sacramental na Eucaristia, hierárquica no Papa e bispos.

Feridas à unidade: cisma, heresia; mas Cristo ora pela unidade, Espírito a realiza.

Santidade: de Cristo Cabeça, Espírito santificador; membros santificados pelos sacramentos, chamados à santidade universal.

Santos canonizados e fiéis defuntos intercedem; Igreja militante, sofredora e triunfante em comunhão.

Para catecúmenos adultos, essas marcas no scrutinium quaresmal unem à unidade eclesial, invocando santidade para purificação sacramental na Vigília Pascal.

\section{A Igreja é Católica e Apostólica (\S 830--856)}

Católica (``universal''): plenitude de Cristo em fé, sacramentos, ministério ordenado; missionária a todos os povos.10 11

Igrejas particulares (dioceses) são católicas unidas ao Sucessor de Pedro; catolicidade histórica em expansão e unidade.

Apostólica: fundada nos Apóstolos, guarda depósito da fé, guiada por bispos sucessores em comunhão petrina.

Sucessão apostólica garante transmissão autêntica do Evangelho e autoridade magisterial.

Apóstolos enviados por Cristo, continuadores de sua missão de pregação e santificação.

Em catequese RICA adulta, a catolicidade e apostolicidade no rito de eleição ancoram fé na tradição viva, motivando compromisso missionário local e universal.

\section{A Missão Apostólica da Igreja (\S 857--870)}

Apóstolos escolhidos por Jesus, testemunhas da Ressurreição, recebem Espírito em Pentecostes para missão universal.

Pedro como pedra, chefe do colégio apostólico; sucessão no Bispo de Roma garante unidade.

Bispos, em comunhão com Papa, sucedem Apóstolos no ensino, santificação, governo.

Presbíteros e diáconos auxiliam; Igreja local (diocese) realiza catolicidade.

Missão: proclamar Evangelho, batizar, formar Povo de Deus até Parusia.

Para adultos no RICA, meditar missão apostólica no tempo pascal da mistagogia impulsiona neófitos ao testemunho, integrando-os na diocese como enviados.

\section{Os Fiéis Leigos e sua Vocação (\S 871--896)}

Fiéis cristãos: batizados, povo de Deus, participam ofícios de Cristo conforme condição.

Igualdade de dignidade, diversidade de funções: hierarquia governa, leigos santificam mundo.

Leigos: secularidade própria, levedura no mundo, testemunhas do Reino pela vida familiar, profissional, social.

Vocação leiga: santificar realidades temporais, exercer paternidade responsável, política justa.

Participação na missão: apostolado hierárquico e leigo, colaboração essencial.

Na iniciação RICA para adultos casados ou profissionais, a vocação leiga no catecumenato desperta apostolado cotidiano, unindo trabalho à liturgia dominical.

\section{A Vida Consagrada (\S 897--913)}

Consagrados: religiosos e seculares, professam conselhos evangélicos (pobreza, castidade, obediência) por voto público.

Testemunham Reino escatológico, seguem Cristo mais intimamente, enriquecem Igreja com carisma.

Vida comum, oração, trabalho; institutos reconhecidos pela Igreja.

Mulheres e homens consagrados: monges, freiras, irmãos, eremitas, virgens consagradas.

Complementam hierarquia e leigos na missão salvífica.

Para catecúmenos adultos discernindo vocações, contemplar consagrados no rito de envio inspira generosidade total, enriquecendo comunidade RICA com diversidade vocacional.

\section{A Vocação e Missão dos Leigos na Igreja e no Mundo (\S 914--945)}

Leigos chamados a santificar mundo desde dentro: matrimônio, família, trabalho, cultura, política.

Discernimento vocacional: graça de estado, direção espiritual.16

Associações leigas: movimentos, grupos para apostolado organizado.

Mulheres leigas: igual dignidade, participação ativa na Igreja e sociedade.

Ministérios leigos: leitores, acólitos, catequistas; serviço eclesial.

Em RICA adulta, essa missão leiga culmina na Crisma, ungindo para transformar realidades seculares com Evangelho, rumo à vida eterna professada.

\section{Conclusão}

Professar a santa Igreja católica revela comunhão trinitária em Povo de Deus una, santa, católica, apostólica. Adultos no RICA, integrem-se como leigos ou consagrados na missão eclesial pascal.